% ============================ CONCLUSION ==============================
\chapter*{Conclusion générale et perspectives}
\addcontentsline{toc}{chapter}{Conclusion générale et perspectives}
\markboth{CONCLUSION GÉNÉRALE}{}

Ce projet de fin d'études a permis de concevoir et de réaliser, au sein de la
\textbf{BTK -- Banque Tuniso-Koweïtienne}, une application web de gestion des
agences bancaires adossée à un volet décisionnel, conduite selon la méthodologie
\textbf{CRISP-DM}. La solution centralise la gestion des agences, des employés, des
clients et des objectifs, assure le \textbf{pointage} des employés, encadre les
opérations sensibles par des circuits de validation et restitue des indicateurs
d'aide à la décision via un tableau de bord et Power BI. Elle intègre en outre une
chaîne \textbf{ETL} et une brique d'\textbf{intelligence artificielle} : la
segmentation des agences par apprentissage non supervisé (clustering).

Ce travail a été l'occasion de mettre en pratique les compétences de la formation
--- modélisation UML, développement Jakarta EE, base de données Oracle, chaîne
décisionnelle (ETL en Python et \emph{reporting} Power BI) et initiation à
l'\textbf{apprentissage automatique} --- au croisement du développement logiciel
et de la Business Intelligence.

Plusieurs \textbf{difficultés} ont jalonné le projet et ont constitué autant
d'occasions d'apprentissage. La \textbf{consolidation} de données dispersées à des
grains hétérogènes a nécessité une chaîne ETL robuste et rejouable. L'articulation
entre le système opérationnel Oracle et l'outil de restitution Power BI a demandé
une attention particulière à la \textbf{modélisation décisionnelle} (schéma en
étoile, mesures) et à la préparation des données. Enfin, la mise au point de la
segmentation a exigé une démarche méthodique de choix et de comparaison de modèles
afin d'en garantir la robustesse. Ces défis ont renforcé une compétence clé :
\textbf{transformer une donnée brute en information utile à la décision}.

Plusieurs perspectives d'évolution peuvent être envisagées :
\begin{itemize}[noitemsep]
  \item le chiffrement (hachage) des mots de passe et le renforcement de la
        sécurité ;
  \item l'\textbf{industrialisation} de la chaîne ETL : ordonnancement automatique
        et historisation incrémentale des indicateurs ;
  \item l'enrichissement par des analyses prédictives (prévision des objectifs) ;
  \item l'exposition d'une API REST pour l'interopérabilité avec d'autres
        systèmes ;
  \item le développement d'une application mobile de consultation des indicateurs.
\end{itemize}
